\chapter{Inteligência Artificial e Metacognição Computacional}

\begin{quote}
\textit{``A metacognição é a inteligência refletindo sobre si mesma — o que torna possível não apenas aprender, mas aprender a aprender.''} --- John Flavell (1979)
\end{quote}

\section{Arquitetura Transformer: O Coração dos LLMs Modernos}

\subsection{Do Token ao Embedding}

O pipeline Transformer começa com a \textbf{tokenização} do texto em subpalavras (Byte-Pair Encoding; Sennrich et al., 2016 \cite{sennrich2016neural}) e a projeção destes tokens em embeddings densos:

\begin{equation}
\mathbf{x} = \text{Embedding}(t) + \text{PositionalEncoding}(p)
\label{eq:transformer_input}
\end{equation}

O \textbf{Positional Encoding} original (Vaswani et al., 2017) usa senos e cossenos:

\begin{equation}
PE_{(pos, 2i)} = \sin\left(\frac{pos}{10000^{2i/d}}\right), \quad PE_{(pos, 2i+1)} = \cos\left(\frac{pos}{10000^{2i/d}}\right)
\label{eq:positional_encoding}
\end{equation}

\subsection{Multi-Head Attention}

O \textbf{MHA} projeta Q, K, V em $h$ cabeças paralelas, permitindo que o modelo ``preste atenção'' em diferentes aspectos simultaneamente:

\begin{equation}
\text{MHA}(Q, K, V) = \text{Concat}(\text{head}_1, \ldots, \text{head}_h) W^O
\label{eq:mha}
\end{equation}
\begin{equation}
\text{head}_i = \text{Attention}(QW_i^Q, KW_i^K, VW_i^V)
\label{eq:mha_head}
\end{equation}

\subsection{Dos Scaling Laws ao Chinchilla}

Kaplan et al. (2020) \cite{kaplan2020scaling} descobriram leis de potência entre loss, parâmetros e dados. Hoffmann et al. (2022) \cite{hoffmann2022training} refinaram com \textbf{Chinchilla}: para um orçamento computacional $C$, o ótimo é:

\begin{equation}
N_{\text{opt}}(C) \propto C^{0.5}, \quad D_{\text{opt}}(C) \propto C^{0.5}
\label{eq:chinchilla}
\end{equation}

Ou seja, parâmetros e tokens de treinamento devem crescer na mesma proporção — ao contrário do GPT-3 (175B parâmetros, sub-treinado).

\section{Reflexion Framework: Agentes que Aprendem com os Próprios Erros}

Shinn et al. (2023) \cite{shinn2023reflexion} introduziram \textbf{Reflexion}: um framework onde agentes mantêm memória episódica de falhas e usam reflexão textual para melhorar. O ciclo é:

\begin{enumerate}
    \item \textbf{Actor}: gera ação/trajetória baseado no estado atual;
    \item \textbf{Evaluator}: avalia o resultado (heurística ou LLM);
    \item \textbf{Self-Reflection}: se falha, gera reflexão textual sobre \textit{por que} falhou;
    \item \textbf{Memory}: armazena reflexão na memória episódica;
    \item \textbf{Repeat}: Actor usa contexto + memória para tentar novamente.
\end{enumerate}

O OpenCode implementa este ciclo via \textbf{Reflexion Engine} (\texttt{mci/reflexion.py}): após cada task, o agente gera reflexão, que é armazenada no MetaBus e usada para melhorar futuras delegações.

\section{Memória Episódica vs Semântica em Agentes Artificiais}

\subsection{A Taxonomia de Tulving Adaptada}

Tulving (1972) \cite{tulving1972episodic} distinguiu memória episódica (eventos pessoais, ``o que, onde, quando'') da memória semântica (fatos, conceitos, conhecimento geral). O OpenCode mapeia:

\begin{itemize}
    \item \textbf{Episódica}: reflexões de execução de tasks — armazenadas em \texttt{MetacognitiveMemory.episodic} (últimos 1000 eventos);
    \item \textbf{Semântica}: lições consolidadas e conhecimento do domínio — armazenadas em \texttt{MetacognitiveMemory.semantic} (dicionário por tópico com \texttt{extract\_lessons()}).
\end{itemize}

\subsection{RAG e Memória Externa}

\textbf{Retrieval-Augmented Generation} (Lewis et al., 2020) \cite{lewis2020retrieval} combina LLMs com busca em bases de conhecimento externas. O OpenCode integra RAG via \texttt{research/llm\_client.py} com Ollama, permitindo que agentes consultem conhecimento externo sem enviar dados para APIs de terceiros.

\section{Global Workspace Theory Implementada em Código}

\subsection{Da Neurociência à Engenharia}

A \textbf{Global Workspace Theory} (GWT; Baars, 1988 \cite{baars1988cognitive}; Dehaene \& Changeux, 2011 \cite{dehaene2011experimental}) postula que a consciência é um ``teatro'' onde informações competem por acesso ao palco central (o \textit{global workspace}) e, uma vez lá, são broadcasted para todos os módulos especializados.

O OpenCode implementa GWT literalmente:

\begin{itemize}
    \item \textbf{Global Workspace} = MetaBus (singleton pub/sub);
    \item \textbf{Módulos especializados} = 128+ agentes (cada um inscrito em tópicos específicos);
    \item \textbf{Competição} = CFP no Blackboard + Attention Router (Trust Score decide vencedor);
    \item \textbf{Broadcast} = MetaBus publica resultados para todos os subscribers.
\end{itemize}

\subsection{A Conexão com o Transformer}

Há um isomorfismo profundo entre a arquitetura Transformer e a GWT:

\begin{itemize}
    \item \textbf{Attention} = Competição por acesso (quais tokens/agentes recebem ``atenção'');
    \item \textbf{Softmax} = Seleção competitiva (winner-take-most, não winner-take-all);
    \item \textbf{Feed-Forward} = Processamento especializado após broadcast.
\end{itemize}

O OpenCode explora este isomorfismo no módulo \texttt{transformer/}, usando camadas Transformer para modelar o fluxo de atenção entre agentes no MetaBus — uma metacognição de segunda ordem.

\begin{thebibliography}{99}
\bibitem{vaswani2017attention} Vaswani, A., et al. (2017). Attention Is All You Need. \textit{NeurIPS 2017}. \doiurl{10.48550/arXiv.1706.03762}
\bibitem{devlin2019bert} Devlin, J., et al. (2019). BERT. \textit{NAACL-HLT 2019}. \doiurl{10.18653/v1/N19-1423}
\bibitem{brown2020language} Brown, T. B., et al. (2020). Language Models are Few-Shot Learners. \textit{NeurIPS 2020}. \doiurl{10.48550/arXiv.2005.14165}
\bibitem{kaplan2020scaling} Kaplan, J., et al. (2020). Scaling Laws for Neural Language Models. \textit{arXiv:2001.08361}. \doiurl{10.48550/arXiv.2001.08361}
\bibitem{hoffmann2022training} Hoffmann, J., et al. (2022). Training Compute-Optimal Large Language Models. \textit{NeurIPS 2022}. \doiurl{10.48550/arXiv.2203.15556}
\bibitem{shinn2023reflexion} Shinn, N., et al. (2023). Reflexion: Language Agents with Verbal RL. \textit{NeurIPS 2023}. \doiurl{10.48550/arXiv.2303.11366}
\bibitem{madaan2023self} Madaan, A., et al. (2023). Self-Refine: Iterative Refinement with Self-Feedback. \textit{NeurIPS 2023}. \doiurl{10.48550/arXiv.2303.17651}
\bibitem{lewis2020retrieval} Lewis, P., et al. (2020). Retrieval-Augmented Generation for Knowledge-Intensive NLP Tasks. \textit{NeurIPS 2020}. \doiurl{10.48550/arXiv.2005.11401}
\bibitem{baars1988cognitive} Baars, B. J. (1988). \textit{A Cognitive Theory of Consciousness}. Cambridge University Press.
\bibitem{dehaene2011experimental} Dehaene, S., \& Changeux, J. P. (2011). Experimental and Theoretical Approaches to Conscious Processing. \textit{Neuron}, 70(2), 200--227. \doiurl{10.1016/j.neuron.2011.03.018}
\bibitem{bengio2023gwt} Bengio, Y. (2023). Global Workspace in AI. \textit{arXiv:2503.03459}.
\bibitem{tulving1972episodic} Tulving, E. (1972). Episodic and Semantic Memory. In \textit{Organization of Memory}. Academic Press.
\bibitem{sennrich2016neural} Sennrich, R., et al. (2016). Neural Machine Translation of Rare Words with Subword Units. \textit{ACL 2016}. \doiurl{10.18653/v1/P16-1162}
\bibitem{flavell1979metacognition} Flavell, J. H. (1979). Metacognition and Cognitive Monitoring. \textit{American Psychologist}, 34(10), 906--911. \doiurl{10.1037/0003-066X.34.10.906}
\end{thebibliography}
